\documentclass{article}

\usepackage[utf8]{inputenc}
\usepackage[T1]{fontenc}
\usepackage{lmodern}
\usepackage{amsmath, amssymb}
\usepackage{microtype}
\usepackage{graphicx}
\usepackage{booktabs}
\usepackage{tabularx}
\usepackage{array}
\usepackage{algorithm}
\usepackage{algpseudocode}
\usepackage{enumitem}
\usepackage{hyperref}

\title{A Framework for Rigorous Evaluation of Model-Agnostic Explainability Methods:\\
Multi-Metric Statistical Benchmarking, Operational Protocol, and Reproducibility}

\author{Paper A Draft}
\date{February 13, 2026}

\begin{document}
\maketitle

\begin{abstract}
Evaluating explainability methods requires more than a single faithfulness proxy. We present a modular benchmarking framework centered on quantitative XAI quality metrics: fidelity, stability, sparsity, computational cost, and faithfulness gap, plus an explicit method for operating the framework end-to-end. On the UCI Adult benchmark, we use a staged protocol (EXP1 calibration/reproducibility and EXP2 comparative/robustness benchmarking); the robustness cohort currently contains 250 of 300 planned configurations (83.3\% coverage). Across complete model-size blocks ($5$ models, $N \in \{50,100,200\}$), Friedman tests indicate significant method differences for fidelity ($\chi^2=42.12$, $p=3.78\times10^{-9}$), stability ($\chi^2=43.88$, $p=1.60\times10^{-9}$), sparsity ($\chi^2=35.64$, $p=8.92\times10^{-8}$), faithfulness gap ($\chi^2=45.00$, $p=9.25\times10^{-10}$), and runtime ($\chi^2=27.72$, $p=4.16\times10^{-6}$). SHAP leads on fidelity/stability, DiCE leads on sparsity, and LIME is generally fastest and most practical outside SVM-KernelSHAP bottlenecks. We release the framework, operation protocol, and artifacts with explicit data-quality caveats for reproducible benchmark use. LLM-based semantic scoring is deferred to a dedicated follow-up study.
\end{abstract}

\section{Introduction}
XAI evaluation remains fragmented across incompatible metrics and ad hoc setups. Most comparisons over-index on fidelity and omit stability, sparsity, and semantic interpretability. This work operationalizes a benchmark-first approach:
\begin{enumerate}[leftmargin=*]
    \item A generic trainer/explainer architecture for model-agnostic evaluation.
    \item A prescriptive operation method that defines how the framework is run, audited, and reported.
    \item Multi-metric quantitative scoring with explicit trade-off analysis.
    \item Reproducibility and statistical rigor suitable for benchmark publication.
\end{enumerate}

\section{Related Work and Framework Design}

\subsection{Literature-Grounded Rationale}
The literature motivating this design supports five requirements:
\begin{enumerate}[leftmargin=*]
    \item Multi-metric evaluation is necessary.
    \item Faithfulness is multi-faceted.
    \item Computational tractability must be reported explicitly.
    \item Benchmark transparency and artifact accounting are first-order requirements.
    \item LLM-judge scoring should be treated as a separate calibrated track.
\end{enumerate}

\subsection{Architecture}
The framework decouples model training, explanation generation, and evaluation through a registry/factory design. Current model coverage includes \texttt{logreg}, \texttt{rf}, \texttt{xgb}, \texttt{svm}, and \texttt{mlp}, with \texttt{lime}, \texttt{shap}, \texttt{anchors}, and \texttt{dice} explainers.

\subsection{Metric Suite (Primary)}
Primary quantitative metrics are:
\begin{itemize}[leftmargin=*]
    \item Fidelity
    \item Stability
    \item Sparsity
    \item Cost
    \item Faithfulness gap
\end{itemize}

\subsection{Paper A Operational Outcome: Framework Operation Method (FOM-7)}
Paper A contributes not only benchmark findings, but also a reusable operating method for end-to-end benchmark execution, auditing, and claim formation.

\section{Methodology}

\subsection{Experimental Cohorts and Scope}
We separate evidence into four cohorts:
\begin{itemize}[leftmargin=*]
    \item \textbf{EXP1 pilot/calibration} (\texttt{exp1\_adult/results})
    \item \textbf{EXP1 reproducibility cohort} (\texttt{exp1\_adult/reproducibility})
    \item \textbf{EXP2 comparative} (\texttt{exp2\_comparative/results})
    \item \textbf{EXP2 robustness cohort} (\texttt{exp2\_scaled/results})
\end{itemize}
Paper A inferential claims are driven by EXP2; EXP1 is retained for calibration and reproducibility reporting.

\subsection{Research Questions and Hypotheses}
\textbf{RQ1}: Do explanation methods differ significantly across complementary technical metrics?\\
\textbf{RQ2}: Are metric trade-offs stable across model families and sample sizes?

For each metric $m \in \{$fidelity, stability, sparsity, faithfulness gap, cost$\}$:
\begin{align}
H_{0,m}&:\text{ all methods have equal median rank across blocks},\\
H_{1,m}&:\text{ at least one method differs.}
\end{align}

\subsection{Dataset, Preprocessing, and Leakage Controls}
The benchmark uses UCI Adult via \path{src/data_loading/adult.py} with:
\begin{itemize}[leftmargin=*]
    \item Stratified train/test split (\texttt{test\_size=0.2}, \texttt{stratify=y})
    \item Seed-controlled deterministic splitting
    \item Numerical pipeline: \texttt{StandardScaler}
    \item Categorical pipeline: \texttt{OneHotEncoder} with \texttt{handle\_unknown='ignore'} and \texttt{sparse\_output=False}
    \item Preprocessor fit only on training split
\end{itemize}

\subsection{Experimental Design and Analysis Units}
\subsubsection{Factorial Grid}
\begin{itemize}[leftmargin=*]
    \item Models: \texttt{logreg}, \texttt{rf}, \texttt{xgb}, \texttt{svm}, \texttt{mlp}
    \item Explainers: \texttt{lime}, \texttt{shap}, \texttt{anchors}, \texttt{dice}
    \item Sample sizes: $N \in \{50,100,200\}$
    \item Seeds: $\{42,123,456,789,999\}$
    \item Planned grid: $5 \times 4 \times 5 \times 3 = 300$ runs
\end{itemize}

\subsubsection{Units for Inference}
\begin{itemize}[leftmargin=*]
    \item Primary unit: run-level metric mean
    \item Block for Friedman: $(\text{model}, N)$
    \item Paired unit for SHAP vs LIME: matched $(\text{model},\text{seed},N)$
\end{itemize}

\subsection{Explainer and Metric Operationalization}
Metrics are computed in \texttt{MetricsEngine} and stored per instance.
\begin{align}
\Delta_k &= \left|f(x)-f\left(x_{\text{mask-top-}k}\right)\right|,\\
\rho &= \mathrm{corr}\left(|w_i|,\left|f(x)-f(x_{-i})\right|\right),\\
S &= \frac{2}{T(T-1)}\sum_{a<b}\cos(e^{(a)},e^{(b)}).
\end{align}
Current implementation maps \texttt{fidelity} to $\rho$ and reports \texttt{faithfulness\_gap} as $\Delta_k$.

\subsection{Artifact Integrity and Coverage Protocol}
Current repository inventory:
\begin{itemize}[leftmargin=*]
    \item \path{experiments/exp1_adult/results}: 29 files (4 core, 24 tuning, 1 diagnostic)
    \item \path{experiments/exp1_adult/reproducibility}: 36 core repeated runs (+1 excluded test run)
    \item \path{experiments/exp2_comparative/results}: 20 files (complete $5\times4$ matrix)
    \item \path{experiments/exp2_scaled/results}: 250 files (of 300 planned)
    \item One malformed JSON: \path{experiments/exp2_scaled/results/}\\
    \path{svm_shap/seed_999/n_200/results.json}
\end{itemize}
Within the 250 robustness files, 233 are analyzable for metric inference, 16 are empty artifacts (no instance evaluations), and 1 is malformed JSON.
Figure~\ref{fig:coverage} makes the completion imbalance explicit at model--explainer granularity.

\begin{figure}[t]
\centering
\includegraphics[width=0.95\linewidth]{../../outputs/paper_a_figures/figures/fig_a2_coverage_heatmap.pdf}
\caption{Coverage map of the EXP2 robustness grid by model and explainer. Cell text reports files present out of 15 expected (\(5\) seeds \(\times\) \(3\) sample sizes), with invalid artifact counts where applicable.}
\label{fig:coverage}
\end{figure}

\subsection{Statistical Analysis Pipeline}
\begin{itemize}[leftmargin=*]
    \item Global method comparison: Friedman test over $(\text{model},N)$ blocks
    \item Post-hoc localization: Nemenyi test
    \item Focused pairwise contrasts: Wilcoxon signed-rank for SHAP vs LIME
    \item Uncertainty reporting: mean, standard deviation, CV, and confidence intervals
\end{itemize}

\subsection{Reproducibility Protocol}
Reproducibility is controlled via fixed seeds, declared configs, structured artifacts, and repeated-run CV reporting.

\subsection{Validity Threats and Mitigations}
\begin{enumerate}[leftmargin=*]
    \item Incomplete factorial coverage
    \item Runtime heterogeneity across explainers
    \item Metric-definition ambiguity
\end{enumerate}

\subsection{Framework Operation Method (FOM-7): Executable Protocol}
\begin{algorithm}[t]
\caption{FOM-7: Framework Operation Method}
\label{alg:fom7}
\begin{algorithmic}[1]
\Require Protocol configuration set $\mathcal{C}$, metric set $\mathcal{M}$, reproducibility seed policy $\mathcal{S}$
\Ensure Claim-ready benchmark report $\mathcal{R}$ with caveat ledger
\State \textbf{Stage S1:} Freeze protocol design and enumerate planned run grid from $\mathcal{C}$
\State \textbf{Gate G1:} Planned run count and all factors are explicit
\State \textbf{Stage S2:} Execute batch runs and persist per-run artifacts and execution manifest
\State \textbf{Gate G2:} Every planned run has terminal state (\texttt{success}/\texttt{failed}/\texttt{skipped})
\State \textbf{Stage S3:} Run integrity audit for missing and malformed artifacts
\State \textbf{Gate G3:} Missing/malformed inventory is explicit and versioned
\State \textbf{Stage S4:} Harmonize outputs and aggregate analysis-ready metric tables
\State \textbf{Gate G4:} All retained runs contain required metrics in $\mathcal{M}$
\State \textbf{Stage S5:} Run inferential pipeline (Friedman, Nemenyi, Wilcoxon as applicable)
\State \textbf{Gate G5:} Inferential outputs are complete for all primary metrics
\State \textbf{Stage S6:} Compute reproducibility statistics on repeated runs
\State \textbf{Gate G6:} Dispersion/CV are reported for each core metric
\State \textbf{Stage S7:} Build claim-ready report with traceability links to artifacts and caveat ledger
\State \textbf{Gate G7:} Every quantitative claim maps to a versioned artifact
\State \Return $\mathcal{R}$
\end{algorithmic}
\end{algorithm}

\begin{figure}[t]
\centering
\includegraphics[width=0.98\linewidth]{../../outputs/paper_a_figures/figures/fig_a1_fom7_flow.pdf}
\caption{Operational view of FOM-7. The workflow is stage-gated: each stage must produce auditable artifacts before progression to the next stage.}
\label{fig:fom7-flow}
\end{figure}
Algorithm~\ref{alg:fom7} defines the formal protocol, and Figure~\ref{fig:fom7-flow} provides its operational view.

\paragraph{Reference implementation command sequence.}
\begin{enumerate}[leftmargin=*]
    \item \texttt{python3 scripts/run\_batch\_experiments.py}\\
    \texttt{--config-dir configs/experiments/exp2\_scaled}\\
    \texttt{--output outputs/batch\_results.csv --parallel}
    \item \texttt{python3 scripts/check\_missing\_results.py}
    \item \texttt{python3 scripts/quick\_summary.py}
    \item \texttt{python3 scripts/full\_summary.py}
    \item \texttt{python3 scripts/run\_reproducibility\_study.py}\\
    \texttt{--config-dir configs/experiments --pattern "exp1\_adult\_*.yaml"}\\
    \texttt{--output-dir experiments/exp1\_adult/reproducibility}
    \item \texttt{python3 scripts/generate\_paper\_a\_figures.py}
\end{enumerate}

\section{Results}

\subsection{Global Trade-offs (EXP2 robustness)}
\begin{figure}[t]
\centering
\includegraphics[width=0.90\linewidth]{../../outputs/paper_a_figures/figures/fig_a3_method_radar.pdf}
\caption{Normalized method profile across core metrics (higher is better after direction-aware normalization). SHAP dominates fidelity/stability, while LIME and DiCE dominate efficiency/sparsity dimensions, respectively.}
\label{fig:method-radar}
\end{figure}
Figure~\ref{fig:method-radar} provides a direction-normalized profile view. Table~\ref{tab:global-tradeoffs} reports absolute means over 15 complete model-size blocks.

\begin{table}[h]
\centering
\caption{Method-level means over complete model-size blocks.}
\begin{tabular}{lrrrrr}
\toprule
Method & Fidelity & Stability & Sparsity & Faithfulness Gap & Time (ms) \\
\midrule
SHAP    & 0.8176 & 0.7377 & 0.2264 & 0.4474 & 685220.23 \\
LIME    & 0.5602 & 0.0144 & 0.0846 & 0.3342 & 3660.68 \\
Anchors & 0.3853 & 0.0006 & 0.0928 & 0.2382 & 25326.92 \\
DiCE    & 0.1716 & 0.3602 & 0.0164 & 0.0988 & 16306.50 \\
\bottomrule
\end{tabular}
\label{tab:global-tradeoffs}
\end{table}

\begin{figure}[t]
\centering
\includegraphics[width=0.88\linewidth]{../../outputs/paper_a_figures/figures/fig_a4_fidelity_cost_tradeoff.pdf}
\caption{Fidelity--cost trade-off frontier using method means over complete blocks (error bars: block-level standard deviation). The x-axis is logarithmic due strong runtime heterogeneity across methods.}
\label{fig:fidelity-cost}
\end{figure}
Figure~\ref{fig:fidelity-cost} isolates the fidelity--cost frontier and its dispersion across complete blocks.

\subsection{Statistical Significance}
\begin{table}[h]
\centering
\caption{Friedman tests across methods.}
\begin{tabular}{lrrl}
\toprule
Metric & Statistic & p-value & Significant \\
\midrule
Fidelity         & 42.12 & 3.78e-09 & Yes \\
Stability        & 43.88 & 1.60e-09 & Yes \\
Sparsity         & 35.64 & 8.92e-08 & Yes \\
Faithfulness Gap & 45.00 & 9.25e-10 & Yes \\
Cost             & 27.72 & 4.16e-06 & Yes \\
\bottomrule
\end{tabular}
\label{tab:friedman}
\end{table}

\subsection{Focused SHAP vs LIME Comparison}
\begin{table}[h]
\centering
\caption{Paired Wilcoxon comparison on 45 matched configurations (\texttt{logreg/rf/xgb}).}
\begin{tabular}{lrrr}
\toprule
Metric & SHAP Mean & LIME Mean & Wilcoxon p-value \\
\midrule
Fidelity         & 0.8112 & 0.5556 & 5.68e-14 \\
Stability        & 0.8013 & 0.0154 & 5.68e-14 \\
Sparsity         & 0.3156 & 0.0845 & 5.68e-14 \\
Faithfulness Gap & 0.3924 & 0.3408 & 5.68e-14 \\
Cost (ms)        & 1258.72 & 210.45 & 2.29e-06 \\
\bottomrule
\end{tabular}
\label{tab:shap-lime}
\end{table}

\subsection{Reproducibility Snapshot (EXP1 repeated runs)}
Fidelity CV $\leq 0.0634$, stability CV $\leq 0.0826$, sparsity CV $\leq 0.0441$, and faithfulness-gap CV $\leq 0.0358$ across baseline configurations. Cost CV reaches 0.225 for SHAP variants.

\begin{figure}[t]
\centering
\includegraphics[width=0.92\linewidth]{../../outputs/paper_a_figures/figures/fig_a5_reproducibility_cv.pdf}
\caption{Reproducibility heatmap for EXP1 repeated runs (9 seeds per core configuration). Lower CV indicates stronger run-to-run stability.}
\label{fig:repro-cv}
\end{figure}
Figure~\ref{fig:repro-cv} highlights that cost variance dominates quality-metric variance in repeated runs.

\section{Validity and Reporting Caveats}
\begin{enumerate}[leftmargin=*]
    \item EXP2 robustness is incomplete (50 missing runs) with one malformed result file.
    \item Runtime comparisons are sensitive to implementation details and hardware.
    \item Semantic evaluation is out of scope for Paper A.
\end{enumerate}

\section{JMLR Track Positioning}
\begin{enumerate}[leftmargin=*]
    \item Reproducible benchmark framework with artifact lineage and failure accounting.
    \item Multi-metric protocol with non-parametric inferential testing.
    \item Reusable operation method (FOM-7) for benchmark execution and audit standardization.
    \item Clearly scoped quantitative benchmark; semantic track deferred.
\end{enumerate}

\section{Conclusion}
The evidence supports a clear trade-off frontier: SHAP leads fidelity/stability, LIME leads runtime practicality, and DiCE leads sparsity. In addition to empirical findings, Paper A contributes FOM-7 as a concrete operating protocol for running, auditing, and reporting XAI benchmark studies.

% Replace with your bibliography when integrating into the final paper package.
% \bibliographystyle{plain}
% \bibliography{references}

\end{document}
